\section{Introduction}
\label{ch:introduction}

\subsection{Background and Context}
\label{sec:background}

South Africa regulates consumer credit tightly. The \ac{NCA} \cite{sa_nca_2005} requires every
registered credit provider to assess whether an applicant can afford a proposed agreement before
granting it, and Regulation 23A of the Affordability Assessment Regulations
\cite{sa_ncr_affordability_2015} prescribes a residual-income test for doing so. Registered lenders
report their exposures to the credit bureaus, so each lender assesses an applicant against a record
of what every other registered lender has advanced.

\ac{BNPL} providers sit outside that regime. Because they levy neither interest nor fees on the
deferred amount and typically rely on late payment penalties, they argue that they do not transact credit within the meaning of the \ac{NCA}, and
so do not register with the \ac{NCR}, do not perform the Regulation 23A assessment, and do not
report what they have lent to any credit bureau \cite{nortonrose_bnpl_sa}. TransUnion, a consumer
credit bureau, confirms the absence of \ac{BNPL} information from South African credit reports
\cite{transunion_cps_sa_2025}. A registered lender running a lawful and diligent affordability
assessment therefore computes residual income from a liability record that is incomplete by
regulatory construction. It cannot see \ac{BNPL} obligations, and neither can any \ac{BNPL} provider
see another's, so a household can hold several concurrent facilities that no party to any of those
agreements has assessed in aggregate.

deHaan et al. \cite{dehaan2024bnpl}, tracking banking
records for 10.6 million United States consumers, find that adopting \ac{BNPL} is followed by higher
overdraft charges, credit-card interest and late fees, concentrated among liquidity-constrained
consumers. The \ac{CFPB} \cite{cfpb2025bnpl} finds that across 2021 and 2022, 63\% of \ac{BNPL}
borrowers held simultaneous loans at some point and 32\% held them across different firms. South
Africa has no equivalent evidence base, and the reporting exemption makes one impossible to
assemble from administrative data, since no party, the regulator included, can measure aggregate
\ac{BNPL} exposure.

\subsection{Research Questions}
\label{sec:rqs}

\begin{quote}
    How does the exemption of \ac{BNPL} from South Africa's credit-reporting and affordability
    regime affect credit distress among \ac{LMI} households, and which platform-level interventions
    reduce it?
\end{quote}

Three subsidiary questions decompose it.

\begin{enumerate}
    \item Can a synthetic household population built from South African survey microdata reproduce
    credit-market behaviour it was not fitted to?
    \item Do households accumulate concurrent \ac{BNPL} facilities that no lender observes in
    aggregate, and does that raise population default?
    \item Which platform-level interventions curb that increase?
\end{enumerate}

The first is a precondition for the other two: a model whose households do not behave like South
African households cannot support a claim about what happens when a new lender class is introduced
to them. It is answered in two parts, by the construction of Section~\ref{ch:population} and by the
baseline of Section~\ref{ch:results}. The second and third are answered in
Section~\ref{ch:results}.